\documentclass[a4paper,10pt]{report}\input{../0.Préambule/Préambule_cours_prof}\begin{document}

\pagestyle{empty}
\newpage
\noindent NOM, Prénom et classe : \dotfill

\section*{Symétrie centrale et fractions \hspace{3cm} Sujet A}
\addcontentsline{toc}{section}{Symétrie centrale et fractions \hspace{3cm} Sujet A}

\noindent \textit{L'usage de la calculatrice est \textbf{interdite}.} \hfill \textit{Durée : 55 minutes}

\medskip
\paragraph{Exercice 1} : \quad Qu'obtient-on quand on fait le symétrique d'un cercle par un point ? \hfill \textbf{2 points}

\paragraph{Exercice 2} : \hfill \textbf{4 points}

\begin{enumerate} 
\item Complète la figure de telle sorte que $O$ en soit le centre de symétrie.
\item En t'aidant du quadrillage, construis le symétrique de la maison par le point $O$
\end{enumerate}

\begin{center}
\includegraphics[scale=0.8]{../40.Image/DS4_1.png}
\end{center}

\paragraph{Exercice 3} : \hfill \textbf{2 points}

\begin{enumerate}
\item Écrit en toute lettre les fractions $\dfrac{2}{5}$ et $\dfrac{4}{3}$
\item Représente les fractions précédentes à l'aides des disques ci dessous. (Tous ne seront pas forcement utile)
\end{enumerate}

\noindent \begin{tikzpicture}[line cap=round,line join=round,>=triangle 45,x=1.0cm,y=1.0cm,scale=0.7]
\clip(3.,2.) rectangle (7.5,6.);
\draw [line width=1.6pt] (5.18,4.02) circle (1.5cm);
\end{tikzpicture}\hspace{0.2cm}
\begin{tikzpicture}[line cap=round,line join=round,>=triangle 45,x=1.0cm,y=1.0cm,scale=0.7]
\clip(3.,2.) rectangle (7.5,6.);
\draw [line width=1.6pt] (5.18,4.02) circle (1.5cm);
\end{tikzpicture}\hspace{0.2cm}
\begin{tikzpicture}[line cap=round,line join=round,>=triangle 45,x=1.0cm,y=1.0cm,scale=0.7]
\clip(3.,2.) rectangle (7.5,6.);
\draw [line width=1.6pt] (5.18,4.02) circle (1.5cm);
\end{tikzpicture}\hspace{0.2cm}
\begin{tikzpicture}[line cap=round,line join=round,>=triangle 45,x=1.0cm,y=1.0cm,scale=0.7]
\clip(3.,2.) rectangle (7.5,6.);
\draw [line width=1.6pt] (5.18,4.02) circle (1.5cm);
\end{tikzpicture}\hspace{0.2cm}
\begin{tikzpicture}[line cap=round,line join=round,>=triangle 45,x=1.0cm,y=1.0cm,scale=0.7]
\clip(3.,2.) rectangle (7.5,6.);
\draw [line width=1.6pt] (5.18,4.02) circle (1.5cm);
\end{tikzpicture}

\vspace{-5mm}
\paragraph{Exercice 4} : \hfill \textbf{3 points}

Sur un collier de $18$ perles, il y a $12$ perles bleues.
\begin{enumerate}
\item Quelle est la proportion de perles bleues sur le collier ?
\item Exprimer cette proportion en pourcentage.
\end{enumerate}

\textbf{Aide :} Les résultats suivant pourront éventuellement être utile :

$12 \times 18 = 216$ \quad\quad $18 \div 12 = 1,5$ \quad\quad $12 \div 18 \approx 0,6667$ \quad\quad $12 \times 30  = 360 $ \quad\quad $12 \div 30 = 0,4$ \quad\quad $30 \div 12 = 2,5$

\vspace{5mm}

\medskip
\noindent \begin{minipage}{7cm}
\paragraph{Exercice 5} : \hfill \textbf{3 points}

Complète le tableau par oui ou non.

\vspace{4mm}
\begin{center}
\renewcommand{\arraystretch}{1.3}
\begin{tabularx}{7cm}{|>{\centering \arraybackslash}m{2.5cm}|*{3}{>{\centering \arraybackslash}X|}}
\hline
Le nombre est-il divisible par \dots & $2$ & $5$ & $9$ \rule[-7pt]{0pt}{18pt}\\\hline
 $19$ & & & \rule[-7pt]{0pt}{18pt}\\\hline
 %$639$ & & & \rule[-7pt]{0pt}{18pt}\\\hline
 $316$ & & & \rule[-7pt]{0pt}{18pt}\\\hline
 $226$ & & & \rule[-7pt]{0pt}{18pt}\\\hline
 %$540$ & & & \rule[-7pt]{0pt}{18pt}\\\hline
 $\np{1785}$ & & & \rule[-7pt]{0pt}{18pt}\\\hline
\end{tabularx}
\end{center}
\end{minipage}
\hspace{10mm}
\begin{minipage}{9cm}

\paragraph{Exercice 6} : \hfill \textbf{3 points}

Pour chaque fraction, coche le (ou les) nombres(s) par le(s)quel(s) elle est simplifiable.

\begin{center}
\begin{tabularx}{\linewidth}{|>{\centering }p{1cm}|*{3}{>{\centering \arraybackslash}X|}}
\hline
 & $\dfrac{42}{60}$ & $\dfrac{30}{60}$ & $\dfrac{189}{99}$  \rule[-7pt]{0pt}{25pt}\\\hline
$2$ & & &  \rule[-7pt]{0pt}{18pt}\\\hline
$3$ & & &  \rule[-7pt]{0pt}{18pt}\\\hline
$5$ & & &  \rule[-7pt]{0pt}{18pt}\\\hline
$9$ & & &  \rule[-7pt]{0pt}{18pt}\\\hline
\end{tabularx}
\end{center}
\end{minipage}

\medskip
\paragraph{Exercice 7} : \quad Simplifie les fractions en utilisant les critères de divisibilité\hfill \textbf{3 points}

\begin{enumerate}
\begin{tabularx}{\linewidth}{*{3}{>{\arraybackslash}X}}
\item $\dfrac{120}{360}$ &
\item $\dfrac{27}{45}$&
\item $\dfrac{135}{27}$ \\
\end{tabularx}
\end{enumerate}

\newpage
\noindent NOM, Prénom et classe : \dotfill

\section*{Symétrie centrale et fractions \hspace{3cm} Sujet B}
\addcontentsline{toc}{section}{Symétrie centrale et fractions \hspace{3cm} Sujet B}

\noindent \textit{L'usage de la calculatrice est \textbf{interdite}.} \hfill \textit{Durée : 55 minutes}

\medskip
\paragraph{Exercice 1} : \quad Qu'obtient-on quand on fait le symétrique d'une droite par un point ? \hfill \textbf{2 points}

\paragraph{Exercice 2} : \hfill \textbf{4 points}

\begin{enumerate} 
\item Complète la figure de telle sorte que $O$ en soit le centre de symétrie.
\item En t'aidant du quadrillage, construis le symétrique de la maison par le point $O$
\end{enumerate}

\begin{center}
\includegraphics[scale=0.8]{../40.Image/DS4_1.png}
\end{center}

\paragraph{Exercice 3} : \hfill \textbf{2 points}

\begin{enumerate}
\item Écrit en toute lettre les fractions $\dfrac{3}{4}$ et $\dfrac{7}{5}$
\item Représente les fractions précédentes à l'aides des disques ci dessous. (Tous ne seront pas forcement utile)
\end{enumerate}

\noindent \begin{tikzpicture}[line cap=round,line join=round,>=triangle 45,x=1.0cm,y=1.0cm,scale=0.7]
\clip(3.,2.) rectangle (7.5,6.);
\draw [line width=1.6pt] (5.18,4.02) circle (1.5cm);
\end{tikzpicture}\hspace{0.2cm}
\begin{tikzpicture}[line cap=round,line join=round,>=triangle 45,x=1.0cm,y=1.0cm,scale=0.7]
\clip(3.,2.) rectangle (7.5,6.);
\draw [line width=1.6pt] (5.18,4.02) circle (1.5cm);
\end{tikzpicture}\hspace{0.2cm}
\begin{tikzpicture}[line cap=round,line join=round,>=triangle 45,x=1.0cm,y=1.0cm,scale=0.7]
\clip(3.,2.) rectangle (7.5,6.);
\draw [line width=1.6pt] (5.18,4.02) circle (1.5cm);
\end{tikzpicture}\hspace{0.2cm}
\begin{tikzpicture}[line cap=round,line join=round,>=triangle 45,x=1.0cm,y=1.0cm,scale=0.7]
\clip(3.,2.) rectangle (7.5,6.);
\draw [line width=1.6pt] (5.18,4.02) circle (1.5cm);
\end{tikzpicture}\hspace{0.2cm}
\begin{tikzpicture}[line cap=round,line join=round,>=triangle 45,x=1.0cm,y=1.0cm,scale=0.7]
\clip(3.,2.) rectangle (7.5,6.);
\draw [line width=1.6pt] (5.18,4.02) circle (1.5cm);
\end{tikzpicture}

\vspace{-5mm}
\paragraph{Exercice 4} : \hfill \textbf{3 points}

Sur un collier de $20$ perles, il y a $12$ perles bleues.
\begin{enumerate}
\item Quelle est la proportion de perles bleues sur le collier ?
\item Exprimer cette proportion en pourcentage.
\end{enumerate}

\textbf{Aide :} Les résultats suivant pourront éventuellement être utile :

$12 \times 20 = 240$ \quad\quad $20 \div 12 \approx 1,667$ \quad\quad $12 \div 20 = 0,6$ \quad\quad $12 \times 32  = 384 $ \quad\quad $12 \div 32 = 0,375$ \quad\quad  $32 \div 12 \approx 2,667$

\vspace{5mm}

\medskip
\noindent \begin{minipage}{7cm}
\paragraph{Exercice 5} : \hfill \textbf{3 points}

Complète le tableau par oui ou non.

\vspace{4mm}
\begin{center}
\renewcommand{\arraystretch}{1.3}
\begin{tabularx}{7cm}{|>{\centering \arraybackslash}m{2.5cm}|*{3}{>{\centering \arraybackslash}X|}}
\hline
Le nombre est-il divisible par \dots & $2$ & $5$ & $9$ \rule[-7pt]{0pt}{18pt}\\\hline
 $23$ & & & \rule[-7pt]{0pt}{18pt}\\\hline
 $639$ & & & \rule[-7pt]{0pt}{18pt}\\\hline
 %$316$ & & & \rule[-7pt]{0pt}{18pt}\\\hline
 %$226$ & & & \rule[-7pt]{0pt}{18pt}\\\hline
 $540$ & & & \rule[-7pt]{0pt}{18pt}\\\hline
 $\np{1765}$ & & & \rule[-7pt]{0pt}{18pt}\\\hline
\end{tabularx}
\end{center}
\end{minipage}
\hspace{10mm}
\begin{minipage}{9cm}

\paragraph{Exercice 6} : \hfill \textbf{3 points}

Pour chaque fraction, coche le (ou les) nombres(s) par le(s)quel(s) elle est simplifiable.

\begin{center}
\begin{tabularx}{\linewidth}{|>{\centering }p{1cm}|*{3}{>{\centering \arraybackslash}X|}}
\hline
 & $\dfrac{44}{60}$ & $\dfrac{90}{180}$ & $\dfrac{189}{93}$  \rule[-7pt]{0pt}{25pt}\\\hline
$2$ & & &  \rule[-7pt]{0pt}{18pt}\\\hline
$3$ & & &  \rule[-7pt]{0pt}{18pt}\\\hline
$5$ & & &  \rule[-7pt]{0pt}{18pt}\\\hline
$9$ & & &  \rule[-7pt]{0pt}{18pt}\\\hline
\end{tabularx}
\end{center}
\end{minipage}

\medskip
\paragraph{Exercice 7} : \quad Simplifie les fractions en utilisant les critères de divisibilité \hfill \textbf{3 points}

\begin{enumerate}
\begin{tabularx}{\linewidth}{*{3}{>{\arraybackslash}X}}
\item $\dfrac{60}{180}$ &
\item $\dfrac{54}{90}$&
\item $\dfrac{135}{27}$ \\
\end{tabularx}
\end{enumerate}
\end{document}
